\documentclass[10pt,a4paper]{amsart}
\usepackage[margin=19mm]{geometry}
\usepackage{amsmath,amssymb,mathtools}
\usepackage[scr=rsfs,cal=euler]{mathalfa}
\usepackage{xfrac}
\usepackage[unicode,hidelinks,bookmarks=false]{hyperref}
\newcommand{\ie}{i.e.\ }
\newcommand{\cf}{cf.\ }
\newcommand{\resp}{respectively\ }
\newcommand{\Subsection}{\subsection}
\providecommand{\nobreakdash}{\nobreak\hbox{-}}
\setlength{\emergencystretch}{2em}
\multlinegap0pt
\usepackage{amssymb}
\usepackage{bm}
\usepackage{xcolor}
\usepackage{mathtools}
\newtheorem{lemma}{Lemma}
\newtheorem{corollary}{Corollary}
\newtheorem{theorem}{Theorem}
\newcommand{\N}{\mathbb{N}}
\newcommand{\nn}{\mathbb{N}}
\newcommand{\qq}{\mathbb{Q}}
\newcommand{\supp}{\text{supp}\,}
\newcommand{\zz}{\mathbb{Z}}
\title{On the set of atoms and strong atoms in additive monoids of cyclic semidomains}
\author{Jiya Dani, Anna Deng, Marly Gotti, Bryan Li, Arav Paladiya, Joseph Vulakh, Jason Zeng}
\date{Source: arXiv:2508.11319v1 (CC BY 4.0)}
\begin{document}
\maketitle

\noindent\emph{Source and licence.} On the set of atoms and strong atoms in additive monoids of cyclic semidomains. Version arXiv:2508.11319v1 (CC BY 4.0), distributed by arXiv under the Creative Commons Attribution 4.0 International licence, http://creativecommons.org/licenses/by/4.0/, as recorded in the arXiv licence metadata for this exact version and shown on the abstract page https://arxiv.org/abs/2508.11319v1. Full licence notice: \texttt{licenses/arxiv-2508.11319-CC-BY-4.0.txt}.\par\smallskip

\noindent\textit{Editorial scope.} Counted unit: the article's theorem thm:4k\_5k\_realize (source0.tex lines 630--632). The article's complete original proof of it is delivered with it (source0.tex lines 634--677, one \texttt{proof} environment of 44 lines). No mathematics is written, repaired or re-derived: the statement and the proof are the article's own lines, in the article's own notation, and no retained line is substituted or re-flowed.  Retained uncounted as the source-local context the proof needs: lines 432--434; lines 448--456; lines 458--460; lines 617--619.  Nothing was omitted from any retained span, so the package carries no omission record.  No retained line contains a citation, so the wrapper carries no bibliography and no bibliography entry.  No retained line contains a figure environment, an image-inclusion command or drawing code, so no image is delivered and the package declares no asset dependency.  Editorial formatting only, in addition: an AMS \texttt{amsart} preamble that restates the article's own declarations and macros used by the retained lines, namely the environments \texttt{lemma}, \texttt{corollary}, \texttt{theorem} on separate counters, together with the standard packages those lines need.  The retained environments print in this document as Lemma 1, Corollary 1, Theorem 1 in order of appearance, whereas the article prints the same items with the numbering of its own full document.  The complete native source of the article as distributed in the arXiv e-print for this exact version is bundled unchanged as \texttt{sources/183316/source0.tex}.
\begin{lemma} \label{lem:strong atom as sum of lower degree terms}
    Let $\alpha$ be an algebraic number with minimal polynomial $m_\alpha(x)$. If~$M_\alpha$ is finitely generated and $1 \le |\mathcal{S}(M_\alpha)| < |\mathcal{A}(M_\alpha)|$, then there exists a polynomial $f(x) \in \qq[x]$ of degree $|\mathcal{S}(M_\alpha)| - \deg m_\alpha(x)$ so that $f(x)m_\alpha(x)$ has negative leading coefficient and all other coefficients nonnegative.
\end{lemma}

\begin{corollary}\label{cor:minimal poly condition for atoms}
    Let $\alpha$ be an algebraic number with minimal polynomial $m_\alpha(x) \in \qq[x]$. Then the following statements hold.
    \begin{enumerate}
        \item If $p(x)$ is a polynomial of minimal degree such that $p(x)m_\alpha(x) = x^n - \sum_{i=0}^{n-1} a_i x^i$ for some $n \in \N$ and coefficients $a_0, \dots, a_{n-1} \in \N_0$, then $|\mathcal{A}(M_\alpha)| = n$.
        \smallskip
        
        \item If $M_\alpha$ is finitely generated and $q(x)$ is a polynomial of minimal degree such that  $q(x)m_\alpha(x) = b_m x^m - \sum_{i=0}^{m-1} b_i x^i$ for some $m \in \N$ and coefficients $b_0, \dots, b_m \in \N_0$, then $|\mathcal{S}(M_\alpha)| = m$.
    \end{enumerate}
\end{corollary}

\begin{corollary}\label{cor:more strong atoms than degree}
    Let $\alpha$ be an algebraic number with minimal polynomial $m_\alpha(x) \in \qq[x]$. If $M_\alpha$ is finitely generated, then $|\mathcal{S}(M_\alpha)| \ge \deg m_\alpha(x)$.
\end{corollary}

\begin{corollary} \label{cor:multiple realizable if eisenstein}
    Let $\alpha$ be an algebraic number with minimal polynomial $m(x) \in \zz[x]$ satisfying the hypothesis of Eisenstein's criterion. Then, for every $k \in \nn$, the polynomial $m(x^k)$ is irreducible in $\qq[x]$ and $f(\beta) = kf(\alpha)$ for any root $\beta$ of $m(x^k)$.
\end{corollary}

\begin{theorem} \label{thm:4k_5k_realize}
    The pair $(4k+c,5k+c)$ is realizable for all $(k,c) \in \nn \times \nn_0$.
\end{theorem}

\begin{proof}
    For the case when $c=0$, we claim that for the positive root $\alpha$ of $m(x) = x^3 - 8x^2 + 4x - 2$, we have $f(\alpha) = (4,5)$. Since $m(x)$ satisfies Eisenstein's criterion, it will follow from Corollary~\ref{cor:multiple realizable if eisenstein} that the pair $(4k,5k)$ is realizable for any $k \in \nn$. To see that $|\mathcal{A}(M_\alpha)| = 5$, note that
    \[
    	(x^2+2x+1)m(x) = x^5 - 6x^4 - 11x^3 - 2x^2 - 2,
    \]
    so $|\mathcal{A}(M_\alpha)|\le 5$. By Corollary~\ref{cor:more strong atoms than degree}, $|\mathcal{A}(M_\alpha)| \ge 3$, and clearly $|\mathcal{A}(M_\alpha)|\ne 3$ as $m(x)$ has multiple negative coefficients. It remains to show that $|\mathcal{A}(M_\alpha)| \ne 4$.
    
    Suppose, for the sake of a contradiction, that there are $4$ atoms. Then there must exist some linear polynomial $f(x) = bx+a \in \zz[x]$ such that $f(x)m(x)$ has a leading coefficient $1$ and all other coefficients non-positive. Multiplying, we get that
    \[
        f(x)m(x) = bx^4 + (a-8b)x^3 + (4b-8a)x^2 + (4a-2b)x - 2a.
    \]
    Note that we must have $b=1$, and also $4b-8a\le 0$ and $4a-2b\le 0$ imply that $b-2a=0$. Thus, clearly $a$ is not an integer, so we reach a contradiction, proving that indeed $|\mathcal{A}(M_\alpha)| = 5$. To show that $|\mathcal{S}(M_\alpha)| = 4$, note that due to Corollary~\ref{cor:more strong atoms than degree}, we have $|\mathcal{S}(M_\alpha)|\ge 3$, and clearly there are more than three strong atoms as $m(x)$ has multiple negative coefficients. Note that $(2x+1)m(x) = 2x^4-15x^3-2$, which satisfies the conditions of Corollary~\ref{cor:minimal poly condition for atoms}, so $|\mathcal{S}(M_\alpha)| = 4$.
    \smallskip
    
    We now consider the case when $c>0$. Fix $(k,c) \in \nn \times \N$, and let us show that the $(4k+c,5k+c)$ belongs to the image of $f$. To do so, first observe that the polynomial
    \[
    	m(x) = x^{3k+c} - 8x^{2k+c} + 4x^{k+c} - 2x^c - 2
    \]
    is irreducible in $\zz[x]$ by Eisenstein's criterion at the ideal $2\zz$ and, therefore, in $\qq[x]$ by virtue of Gauss's lemma. In addition, $m(x)$ has a positive root $\alpha$ as $m(0) < 0$. We proceed to show that $M_\alpha$ is an atomic monoid with exactly $5k+c$ atoms and $4k+c$ strong atoms. To see that there are at most $5k+c$ atoms, consider the polynomial $f(x) = 1 + 2x^k + x^{2k}$, so that
    \[
        f(x)m(x)=-2-2x^c-4x^k-2x^{2k}-2x^{2k+c}-11x^{3k+c}-6x^{4k+c}+x^{5k+c}.
    \]
    Similarly, there are at most $4k+c$ strong atoms since multiplying $m(x)$ by $1+2x^k$ gives
    \[
    	(1+2x^{k}) m(x) = -2-2x^c-4x^k-15x^{3k+c}+2x^{4k+c}.
    \]
    Note that there are no issues in both cases when $c$ is a multiple of $k$, as all non-leading coefficients will remain negative. Now suppose for the sake of contradiction that $\alpha^{4k+c-1}$ is not a strong atom. Then there exists some $q(x)$ so that $m(x)q(x)$ has all non-positive coefficients except for the coefficient of $x^{4k+c-1}$. We can assume $\deg q(x)=k-1$ by Lemma~\ref{lem:strong atom as sum of lower degree terms}, so write $q(x)=\sum_{d=0}^{k-1}a_dx^d$. Then
    \[
    	 m(x) a_dx^d = -2a_dx^d-2a_dx^{d+c}+4a_dx^{k+c+d}-8a_dx^{2k+c+d}+a_dx^{3k+c+d}.
    \]
    Observe that for any $d<k-1$, the set $\supp(a_{k-1}x^{k-1}m(x))\,\cap\,\supp(a_{d}x^{d}m(x))$ contains only numbers less than $k+c+d$. Therefore the coefficients of the $x^{2k+c-1}$ and $x^{3k+c-1}$ terms are given by $4a_{k-1}$, and $-8a_{k-1}$, respectively, from which it follows that $a_{k-1} = 0$. However, this contradicts that $\deg q(x) = k-1$.
    
    Similarly, for the atoms, suppose for the sake of a contradiction that $\alpha^{5k+c-1}$ is not an atom. Then there exists some $q(x)=\sum_{d=0}^{k-1}q_d(x)$ where $q_d(x)=a_dx^d+b_dx^{d+k}$ so that $m(x)q(x)$ has all non-positive coefficients except for the term $x^{5k+c-1}$, which has coefficient~$1$. We have
    \begin{align*}
    		m(x)q_d(x) &= -2a_dx^d-2a_dx^{d+c}-2b_dx^{d+k}+(4a_d-2b_d)x^{k+c+d} \\
    							&+(4b_d-8a_d)x^{2k+c+d}+(a_d-8b_d)x^{3k+c+d}+b_dx^{4k+c+d}.
    \end{align*}
	Note that the set $\supp(q_d(x)m(x))\,\cap\,\supp(q_{d'}(x)m(x))$ contains only numbers less than $2k+c+d'$  for each pair $(d',d)$ with $d'<d<k$. Therefore the coefficients of the terms $x^{2k+c+d}, x^{3k+c+d}$, and $x^{4k+c+d}$ in $m(x)q(x)$ are given by $4a_d-8b_d$, $a_d-8b_d$, and $b_d$, respectively. Hence, when $d<k-1$, the inequalities $b_d \le 0$ and $a_d \le 8b_d$ hold, and so it follows that $a_d \le 0$. Now, consider the coefficient of term $x^d$ for $d<k-1$ in $m(x)q(x)$. Only $q_d(x)$ and $q_{d-c}(x)$ can contribute and so the coefficient of the $x^d$-term of the polynomial $q(x)m(x)$ is $-2a_d - 2a_{d-c}$, which results from writing the $x^d$-term as $-2a_dx^d - 2a_{d-c}x^{(d-c) + c}$ (set $a_{d-c}$ to be $0$ when $d-c < 0$). Suppose that $a_{d-c}=0$ for some $d$, and notice that because the coefficient in $m(x)q(x)$ of the term $x^d$ cannot be positive, $a_d \le 0$, and so $a_d = 0$. Since $a_{d-c}=0$ for $d<c$, it follows by induction that $a_d = 0$ for each $d<k-1$. Additionally, from the inequalities $b_d\le 0$ and $8b_d \ge a_d$ we obtain that $b_d = 0$. Thus, $q(x) = q_{k-1}(x) = a_{k-1}x^{k-1}+x^{2k-1}$. Now multiplying $q(x)$ by $m(x)$ gives 
    \begin{align*}
    	m(x)q(x) &=-2a_{k-1}x^{k-1}-2a_{k-1}x^{{k-1}+c}-2x^{{2k-1}}+(4a_{k-1}-2)x^{{2k+c}-1} \\
    					&+(4-8a_{k-1})x^{3k+c-1}+(a_{k-1}-8)x^{4k+c-1}+x^{5k+c-1}.
    \end{align*}
	We get $a_{k-1} \ge 0$ from the first coefficient, and $a_{k-1}\le\frac12$ from the coefficient of the term $x^{2k+c-1}$. As a result, $a_{k-1} = 0$, which contradicts $4 - 8a_{k-1} \le 0$ from the coefficient of the term $x^{3k+c-1}$.
\end{proof}

\end{document}
